\documentclass[10pt,a4paper]{amsart}
\usepackage[margin=19mm]{geometry}
\usepackage{amsmath,amssymb,mathtools}
\usepackage[scr=rsfs,cal=euler]{mathalfa}
\usepackage{xfrac}
\usepackage[unicode,hidelinks,bookmarks=false]{hyperref}
\newcommand{\ie}{i.e.\ }
\newcommand{\cf}{cf.\ }
\newcommand{\resp}{respectively\ }
\newcommand{\Subsection}{\subsection}
\providecommand{\nobreakdash}{\nobreak\hbox{-}}
\setlength{\emergencystretch}{2em}
\multlinegap0pt
\usepackage[shortlabels]{enumitem}
\usepackage{amssymb}
\usepackage{mathtools}
\newtheorem{theorem}{Theorem}
\newtheorem{lemma}{Lemma}
\newcommand{\C}{\mathbb{C}}
\newcommand{\We}{W_e}
\DeclareMathOperator{\codim}{codim}
\DeclareMathOperator{\conv}{conv}
\newcommand{\e}{\varepsilon}
\newcommand{\ip}[2]{\langle #1,#2\rangle}
\newcommand{\wto}{\xrightarrow{\mathrm w}}
\title{Numerical and essential numerical ranges on $\ell_p$}
\author{Vladimir M\"uller, Yuri Tomilov}
\date{Source: arXiv:2606.22693v1 (CC BY 4.0)}
\begin{document}
\maketitle

\noindent\emph{Source and licence.} Numerical and essential numerical ranges on $\ell_p$. Version arXiv:2606.22693v1 (CC BY 4.0), distributed by arXiv under the Creative Commons Attribution 4.0 International licence, http://creativecommons.org/licenses/by/4.0/, as recorded in the arXiv licence metadata for this exact version and shown on the abstract page https://arxiv.org/abs/2606.22693v1. Full licence notice: \texttt{licenses/arxiv-2606.22693-CC-BY-4.0.txt}.\par\smallskip

\noindent\textit{Editorial scope.} Counted unit: the article's theorem thm:convhull-sigma-pie (source0.tex lines 634--639). The article's complete original proof of it is delivered with it (source0.tex lines 641--805, one \texttt{proof} environment of 165 lines). No mathematics is written, repaired or re-derived: the statement and the proof are the article's own lines, in the article's own notation, and no retained line is substituted or re-flowed.  Retained uncounted as the source-local context the proof needs: lines 472--494; lines 592--596; lines 605--612; lines 614--628.  Nothing was omitted from any retained span, so the package carries no omission record.  No retained line contains a citation, so the wrapper carries no bibliography and no bibliography entry.  No retained line contains a figure environment, an image-inclusion command or drawing code, so no image is delivered and the package declares no asset dependency.  Editorial formatting only, in addition: an AMS \texttt{amsart} preamble that restates the article's own declarations and macros used by the retained lines, namely the environments \texttt{theorem}, \texttt{lemma} on separate counters, together with the standard packages those lines need.  The retained environments print in this document as Theorem 1, Lemma 1 in order of appearance, whereas the article prints the same items with the numbering of its own full document.  The complete native source of the article as distributed in the arXiv e-print for this exact version is bundled unchanged as \texttt{sources/203374/source0.tex}.
\begin{theorem}\label{thm:characterization-We}
%Let $X$ be a Banach space,
 Let $T\in \mathcal L(X)$, and let $\lambda\in\C$. Then the following are equivalent.
\begin{enumerate}[label=\textup{(\roman*)}]
\item $\lambda\in\We(T)$.
\item There exists a net $(x_\alpha,x^*_\alpha)\in X\times X^*$ such that
\[
x_\alpha\wto 0,
\quad
\|x_\alpha\|\to 1,
\quad
\|x_\alpha^*\|\to 1,
\quad
\ip{x_\alpha}{x_\alpha^*}\to 1,
\quad
\ip{Tx_\alpha}{x_\alpha^*}\to \lambda.
\]
\item For every subspace $M\subset X$ with $\codim M<\infty$, and for every $\e>0$, there exists $(x,x^*)\in\Pi(X)$ with $x\in M$ and
\[
|\ip{Tx}{x^*}-\lambda|<\e.
\]
\end{enumerate}
\end{theorem}

\begin{equation}\label{eq:conv-sigmae-sigmapie}
        \operatorname{conv}\sigma_e(T)
        =
        \operatorname{conv}\sigma_{\pi e}(T).
\end{equation}

\begin{lemma}\label{lem:separation-finite-codim}
Let $N\subset X$ be a finite-dimensional subspace and let $0<\delta<1$. Then there exists a finite-codimensional subspace $M\subset X$ such that
\[
\|u+v\|\ge (1-\delta)\|u\|
\qquad (u\in N,\ v\in M).
\]
%In particular, $F\cap N=\{0\}$.
\end{lemma}

\begin{lemma}\label{lem:zenger}
Let $x_1,\dots,x_m$ be linearly independent vectors in $X,$
%a complex Banach space $X$, 
and let $a_1,\dots,a_m>0$ satisfy $\sum_{k=1}^m a_k=1$. Then there exist complex numbers $t_1,\dots,t_m$, a vector
\[
x=\sum_{k=1}^m t_kx_k\in X,
\]
and a functional $f\in X^*$ such that
\[
\|x\|=\|f\|=\ip{x}{f}=1,
\qquad
\ip{t_kx_k}{f}=a_k
\quad (1\le k\le m).
\]
\end{lemma}

\begin{theorem}\label{thm:convhull-sigma-pie}
Let $X$ be a complex Banach space and let $T\in \mathcal L(X)$. Then
\[
\conv\sigma_{e}(T)\subset \We(T).
\]
\end{theorem}

\begin{proof}
By \eqref{eq:conv-sigmae-sigmapie}, it is enough to show that
\[
        \operatorname{conv}\sigma_{\pi e}(T)\subset W_e(T).
\]
Let
\[
        \lambda=\sum_{k=1}^m a_k\lambda_k,
        \qquad
        a_k>0,\quad \sum_{k=1}^m a_k=1,\quad
        \lambda_k\in\sigma_{\pi e}(T),
\]
and let \(M\subset X\) be a finite-codimensional subspace. We prove that,
for every \(\varepsilon>0\), there exists \((x,x^*)\in\Pi(X)\) such that
\(x\in M\) and
\[
        |\langle Tx,x^*\rangle-\lambda|<\varepsilon .
\]
The conclusion will then follow from
Theorem~\ref{thm:characterization-We}.

Choose \(0<\delta<1/2\), and put
\[
        C:=(1-\delta)^{-1},
        \qquad
        \eta:=\frac{\varepsilon}{2mC}.
\]
We construct unit vectors \(x_1,\ldots,x_m\in M\). Set \(L_0:=M\). Since
\(\lambda_1\in\sigma_{\pi e}(T)\), choose \(x_1\in L_0\) such that
\[
        \|x_1\|=1,
        \qquad
        \|(T-\lambda_1I)x_1\|<\eta .
\]
Suppose that \(x_1,\ldots,x_r\) have been chosen, where \(1\le r<m\), and
let
\[
        F_r:=\operatorname{span}\{x_1,\ldots,x_r\}.
\]
By Lemma~\ref{lem:separation-finite-codim}, there is a finite-codimensional
subspace \(N_r\subset X\) such that
\[
        \|u+v\|\ge (1-\delta)\|u\|
        \qquad (u\in F_r,\ v\in N_r).
\]
In particular, \(F_r\cap N_r=\{0\}\). Define
\[
        L_r:=M\cap N_1\cap\cdots\cap N_r .
\]
Then \(L_r\) is finite-codimensional. Since
\(\lambda_{r+1}\in\sigma_{\pi e}(T)\), choose \(x_{r+1}\in L_r\) such that
\[
        \|x_{r+1}\|=1,
        \qquad
        \|(T-\lambda_{r+1}I)x_{r+1}\|<\eta .
\]
This completes the recursive construction.

Since \(x_{r+1}\in N_r\) and \(F_r\cap N_r=\{0\}\) for \(1\le r<m\), the
vectors \(x_1,\ldots,x_m\) are linearly independent. Let
\[
        Y:=\operatorname{span}\{x_1,\ldots,x_m\}.
\]
For \(0\le r\le m\), let \(P_r:Y\to Y\) be the projection onto
\(\operatorname{span}\{x_1,\ldots,x_r\}\) along
\(\operatorname{span}\{x_{r+1},\ldots,x_m\}\), with \(P_0=0\) and \(P_m=I_Y\).
We claim that
\[
        \|P_r\|\le C \qquad (0\le r\le m).
\]
For \(r=0\) and \(r=m\) this is clear. Let \(1\le r<m\), and write
\[
        y=\sum_{k=1}^m \beta_kx_k\in Y
\]
with complex scalars \(\beta_k\). Then
\[
        P_ry\in F_r,
        \qquad
        y-P_ry=\sum_{k=r+1}^m \beta_kx_k\in N_r,
\]
because \(x_k\in L_{k-1}\subset N_r\) for \(k>r\). Hence
\[
        \|y\|\ge (1-\delta)\|P_ry\|,
\]
and the claim follows.

Let \(\varphi_1,\ldots,\varphi_m\in Y^*\) be the coordinate functionals
associated with the basis \(x_1,\ldots,x_m\). Since
\[
        (P_k-P_{k-1})y=\varphi_k(y)x_k,
        \qquad \|x_k\|=1,
\]
we have
\[
        |\varphi_k(y)|
        \le
        \bigl(\|P_k\|+\|P_{k-1}\|\bigr)\|y\|
        \le 2C\|y\|
        \qquad (1\le k\le m),
\]
and therefore
\[
        \|\varphi_k\|\le 2C
        \qquad (1\le k\le m).
\]

By Lemma~\ref{lem:zenger}, there exist complex scalars \(t_1,\ldots,t_m\), a
vector
\[
        x:=\sum_{k=1}^m t_kx_k\in Y,
\]
and a functional \(g\in Y^*\) such that
\[
        \|x\|=\|g\|=\langle x,g\rangle=1,
        \qquad
        \langle t_kx_k,g\rangle=a_k
        \quad (1\le k\le m).
\]
Let \(x^*\in X^*\) be a Hahn--Banach extension of \(g\) with
\(\|x^*\|=1\). Then
\[
        (x,x^*)\in\Pi(X),
        \qquad
        x\in Y\subset M.
\]
Moreover, since \(t_k=\varphi_k(x)\),
\[
        |t_k|\le 2C
        \qquad (1\le k\le m).
\]
Using
\[
        \langle t_kx_k,x^*\rangle
        =
        \langle t_kx_k,g\rangle
        =
        a_k
        \qquad (1\le k\le m),
\]
we obtain
\[
\begin{aligned}
        \langle Tx,x^*\rangle-\lambda
        &=
        \sum_{k=1}^m
        \langle t_k(T-\lambda_kI)x_k,x^*\rangle .
\end{aligned}
\]
Consequently,
\[
\begin{aligned}
        |\langle Tx,x^*\rangle-\lambda|
        &\le
        \sum_{k=1}^m
        |t_k|\,\|(T-\lambda_kI)x_k\|\,\|x^*\|  \\
        &\le
        \sum_{k=1}^m 2C\eta
        =
        \varepsilon .
\end{aligned}
\]
Since \(M\) and \(\varepsilon\) were arbitrary, 
 Theorem~\ref{thm:characterization-We} gives
\(\lambda\in W_e(T)\).
\end{proof}

\end{document}
